\par\vspace{5pt}\noindent\begin{minipage}{\linewidth}
\reghead{IFRL (and IERL, WAITL)}{0x06 (0x22, 0x20)}{W1C / R/W}{\texttt{0x00}}
\vspace{2pt}
{\centering\footnotesize\setlength{\tabcolsep}{2pt}%
\begin{tabular}{@{}|>{\centering\arraybackslash}p{12mm}|*{8}{>{\centering\arraybackslash}p{16.9mm}|}@{}}
\hline
\rowcolor{fillD}\textbf{Bit} & \textbf{7} & \textbf{6} & \textbf{5} & \textbf{4} & \textbf{3} & \textbf{2} & \textbf{1} & \textbf{0} \\ \hline
\textbf{Field} & \texttt{DLY} & \texttt{TRIM} & \texttt{MBX} & \texttt{OVF} & \texttt{OCB} & \texttt{OCA} & \texttt{CAPT} & \texttt{PC} \\ \hline
\end{tabular}\par}
\vspace{2pt}
{\small\begin{tabularx}{\linewidth}{@{}llL@{}}
\toprule
Bits & Field & Description \\
\midrule
\texttt{7} & \reg{DLY} & Shared time base reached the reading core's wake-up time (per core). Cleared by writing a new wake-up time to \reg{DLYL/H} \\
\texttt{6} & \reg{TRIM} & Trim counter completed a measurement inside the target window. Latched, W1C \\
\texttt{5} & \reg{MBX} & Host mailbox holds an unread byte from the host. Level, follows \reg{MBSTAT.RXF} \\
\texttt{4} & \reg{OVF} & 16-bit timer overflow. Latched, W1C \\
\texttt{3:2} & \reg{OCB, OCA} & 16-bit timer compare B and compare A match. Latched, W1C \\
\texttt{1} & \reg{CAPT} & 16-bit timer input capture. Latched, W1C \\
\texttt{0} & \reg{PC} & Pin change on an unmasked bidirectional pin (\reg{PCMSK}). Latched, W1C \\
\bottomrule
\end{tabularx}\par}
\noindent{\footnotesize\reg{IERL} (\texttt{0x22}) enables interrupts and \reg{WAITL} (\texttt{0x20}) sets the wait mask for the same bit positions (one register per core); writing \reg{WAITL} starts the wait.}
\end{minipage}\par

\par\vspace{5pt}\noindent\begin{minipage}{\linewidth}
\reghead{IFRH (and IERH, WAITH)}{0x07 (0x23, 0x21)}{mixed / R/W}{\texttt{0x00}}
\vspace{2pt}
{\centering\footnotesize\setlength{\tabcolsep}{2pt}%
\begin{tabular}{@{}|>{\centering\arraybackslash}p{12mm}|*{8}{>{\centering\arraybackslash}p{16.9mm}|}@{}}
\hline
\rowcolor{fillD}\textbf{Bit} & \textbf{7} & \textbf{6} & \textbf{5} & \textbf{4} & \textbf{3} & \textbf{2} & \textbf{1} & \textbf{0} \\ \hline
\textbf{Field} & \texttt{SW3} & \texttt{SW2} & \texttt{SW1} & \texttt{SW0} & \texttt{LT1} & \texttt{LR1} & \texttt{LT0} & \texttt{LR0} \\ \hline
\end{tabular}\par}
\vspace{2pt}
{\small\begin{tabularx}{\linewidth}{@{}llL@{}}
\toprule
Bits & Field & Description \\
\midrule
\texttt{7:4} & \reg{SW3..SW0} & Software flags shared by both cores (R/W; set with \texttt{SBI}, cleared with \texttt{CBI}) \\
\texttt{3:0} & \reg{LT1..LR0} & Line-cell event flags (a status bit selected in \reg{LCEVM} is set). Level; visible only to the pair's owner core (\reg{LCOWN}) \\
\bottomrule
\end{tabularx}\par}
\noindent{\footnotesize\reg{IERH} (\texttt{0x23}) and \reg{WAITH} (\texttt{0x21}) use the same bit positions per core.}
\end{minipage}\par

\par\vspace{5pt}\noindent\begin{minipage}{\linewidth}
\reghead{DLYL, DLYH}{0x24, 0x25}{R / W (per core)}{\texttt{0x00}}
\vspace{2pt}
{\small\noindent Read returns the shared 16-bit free-running peripheral time base. Write sets this core's wake-up time and clears \reg{IFRL.DLY}\par}
\end{minipage}\par